\ifdefined\gkver\else\def\gkver{student}\fi
\documentclass[\gkver]{gaokaozhenti}
\begin{document}

\chapter{2012年四川理综}
%% paperType: 理综物理部分

\begin{enumerate}
\item
%% number: 14
%% paperName: 2012年四川理综第 14 题 6 分
%% typeId: 1
%% type: 单选题
%% chapter: 气体性质
%% point: 物体的内能
%% method: 无
%% score: 6
%% degree: 650
%% duplicateId: 0
%% body:
物体由大量分子组成，下列说法正确的是\xzanswer{C}
\fourchoices[answer=C]
{分子热运动越剧烈，物体内每个分子的动能越大}
{分子间引力总是随着分子间的距离减小而减小}
{物体的内能跟物体的温度和体积有关}
{只有外界对物体做功才能增加物体的内能}
%% answer: C
%% memo:
\memoanswer{
分子热运动越剧烈，分子的平均动能越大，但并非每个分子的动能都越大，A错误；

分子间引力随分子间的距离减小而增大，B错误；

物体的内能与分子动能和分子势能有关，即与温度和体积有关，C正确；

通过热传递也可以增加物体的内能，D错误。
}

\item
%% number: 15
%% paperName: 2012年四川理综第 15 题 6 分
%% typeId: 1
%% type: 单选题
%% chapter: 曲线运动
%% point: 人造地球卫星
%% method: 无
%% score: 6
%% degree: 650
%% duplicateId: 0
%% body:
今年4月30日，西昌卫星发射中心的中圆轨道卫星，其轨道半径为$2.8\times10^{7}\Um$。它与另一颗同质量的同步轨道卫星（轨道半径为$4.2\times10^{7}\Um$）相比\xzanswer{B}
\fourchoices[answer=B]
{向心力较小}
{动能较大}
{发射速度都是第一宇宙速度}
{角速度较小}
%% answer: B
%% memo:
\memoanswer{
由题目所给的数据知，中圆轨道卫星的轨道半径小于同步轨道卫星的半径，根据万有引力定律，向心力由万有引力提供，即 $F_{\text{向}}=G\frac{Mm}{r^{2}}$，故中圆轨道卫星的向心力较大，A错误；由 $G\frac{Mm}{r^{2}}=m\frac{v^{2}}{r}$，得 $v=\sqrt{\frac{GM}{r}}$，故中圆轨道卫星的线速度较大，动能较大，B正确；由于这两个卫星都不是近地卫星，所以发射速度均大于第一宇宙速度，C错误；由 $\omega=\frac{v}{r}$ 知中圆轨道卫星的角速度较大，D错误。
}

\item
%% number: 16
%% paperName: 2012年四川理综第 16 题 6 分
%% typeId: 1
%% type: 单选题
%% chapter: 电磁感应
%% point: 电磁感应现象
%% method: 无
%% score: 6
%% degree: 650
%% duplicateId: 0
%% body:
如图所示，在铁芯P上绕着两个线圈a和b，则\xzanswer{D}
\onepicture{figs/16.png}
\fourchoices[answer=D]
{线圈a输入正弦交变电流，线圈b可输入恒定电流}
{线圈a输入恒定电流，穿过线圈b的磁通量一定为零}
{线圈b输出的交变电流不对线圈a的磁场造成影响}
{线圈a的磁场变化时，线圈b中一定有电场}
%% answer: D
%% memo:
\memoanswer{
当线圈a输入正弦交变电流时，线圈b输出同频率的正弦交变电流，A错误；当线圈a输入恒定电流时，线圈a产生稳定的磁场，通过线圈b的磁通量不变，但不是零，B错误；由于互感，每个线圈的交变电流都对另外一个线圈的磁场产生影响，C错误；根据麦克斯韦电磁场理论，变化的磁场一定产生电场，D正确。
}

\item
%% number: 17
%% paperName: 2012年四川理综第 17 题 6 分
%% typeId: 1
%% type: 单选题
%% chapter: 原子和原子核
%% point: 结合能
%% method: 无
%% score: 6
%% degree: 650
%% duplicateId: 0
%% body:
如图为氢原子能级示意图的一部分，则氢原子\xzanswer{A}
\onepicture{figs/17.png}
\fourchoices[answer=A]
{从$n=4$能级跃迁到$n=3$能级比从$n=3$能级跃迁到$n=2$能级辐射出电磁波的波长长}
{从$n=5$能级跃迁到$n=1$能级比从$n=5$能级跃迁到$n=4$能级辐射出电磁波的速度大}
{处于不同能级时，核外电子在各处出现的概率是一样的}
{从高能级向低能级跃迁时，氢原子核一定向外放出能量}
%% answer: A
%% memo:
\memoanswer{
由图知，$n=4$能级与$n=3$能级间的能量差小于$n=3$能级与$n=2$能级间的能量差，根据$\Delta E=h\nu=h\frac{c}{\lambda}$，A正确；真空中所有电磁波的速度都相同，B错误；处于不同能级时，氢原子的核外电子在各处出现的概率是不一样的，C错误；从高能级向低能级跃迁时，氢原子向外放出能量，但氢原子核不发生变化，D错误。
}

\item
%% number: 18
%% paperName: 2012年四川理综第 18 题 6 分
%% typeId: 1
%% type: 单选题
%% chapter: 光学
%% point: 光的折射
%% method: 无
%% score: 6
%% degree: 650
%% duplicateId: 0
%% body:
a、b两种单色光组成的光束从介质进入空气时，其折射光束如图所示。用a、b两束光\xzanswer{C}
\onepicture{figs/18.png}
\fourchoices[answer=C]
{先后照射双缝干涉实验装置，在缝后屏上都能出现干涉条纹，由此确定光是横波}
{先后照射某金属，a光照射时恰能逸出光电子，b光照射时也能逸出光电子}
{从同一介质以相同方向射向空气，其界面为平面，若b光不能进入空气，则a光也不能进入空气}
{从同一介质以相同方向射向空气，其界面为平面，a光的反射角比b光的反射角大}
%% answer: C
%% memo:
\memoanswer{
横波和纵波都能发生干涉现象，A错误；由图知，a光的折射率较大，则a光的频率较大，所以若a光照射某金属恰能逸出光电子，则b光照射该金属一定不能逸出光电子，B错误；a光的折射率较大，则a光的临界角较小，以相同入射角从同一个介质射入空气时，若b光发生全反射，则a光一定发生全反射，C正确；根据光的反射定律，反射角等于入射角，D错误。
}

\item
%% number: 19
%% paperName: 2012年四川理综第 19 题 6 分
%% typeId: 1
%% type: 单选题
%% chapter: 机械波
%% point: 波的图象
%% method: 无
%% score: 6
%% degree: 450
%% duplicateId: 0
%% body:
在$xOy$平面内有一列沿$x$轴正方向传播的简谐横波，波速为$2\Ums$，振幅为$A$。M、N是平衡位置相距$2\Um$的两质点，如图所示。在$t=0$时，M通过其平衡位置沿$y$轴正方向运动，N位于其平衡位置上方最大位移处。已知该波的周期大于$1\Us$。则\xzanswer{D}
\onepicture{figs/19.png}
\fourchoices[answer=D]
{该波的周期为$\frac{5}{3}\Us$}
{在$t=\frac{1}{3}\Us$时，N的速度一定为$2\Ums$}
{从$t=0$到$t=1\Us$，M向右移动了$2\Um$}
{从$t=\frac{1}{3}\Us$到$t=\frac{2}{3}\Us$，M的动能逐渐增大}
%% answer: D
%% memo:
\memoanswer{
设波长为$\lambda$，根据题意，$(n+\frac{3}{4})\lambda=2\Um$，结合$v=\frac{\lambda}{T}$，得$T=\frac{1}{n+\frac{3}{4}}\Us>1\Us$，所以$n=0$，$T=\frac{4}{3}\Us$，A错误；

质点振动速度与波速无必然联系，B错误；

波传播过程中，质点并不随波迁移，C错误；

从$\frac{1}{3}\Us$到$\frac{2}{3}\Us$，质点M从波峰向平衡位置运动，速度变大，动能变大，D正确。
}

\item
%% number: 20
%% paperName: 2012年四川理综第 20 题 6 分
%% typeId: 2
%% type: 多选题
%% chapter: 电磁感应
%% point: 电磁感应电路分析
%% method: 无
%% score: 6
%% degree: 450
%% duplicateId: 0
%% body:
半径为$a$右端开小口的导体圆环和长为$2a$的导体直杆，单位长度电阻均为$R_{0}$。圆环水平固定放置，整个内部区域分布着竖直向下的匀强磁场，磁感应强度为$B$。杆在圆环上以速度$v$平行于直径CD向右做匀速直线运动，杆始终有两点与圆环良好接触，从圆环中心O开始，杆的位置由$\theta$确定，如图所示。则\xzanswer{AD}
\onepicture{figs/20.png}
\fourchoices[answer=AD]
{$\theta=0$时，杆产生的电动势为$2Bav$}
{$\theta=\frac{\pi}{3}$时，杆产生的电动势为$\sqrt{3}Bav$}
{$\theta=0$时，杆受的安培力大小为$\frac{3B^{2}av}{(\pi+2)R_{0}}$}
{$\theta=\frac{\pi}{3}$时，杆受的安培力大小为$\frac{3B^{2}av}{(5\pi+3)R_{0}}$}
%% answer: AD
%% memo:
\memoanswer{
当 $\theta=0$ 时，杆在圆心位置，切割磁感线的有效长度等于圆环直径，杆产生的感应电动势为 $E=2Bav$，A正确；当 $\theta=\frac{\pi}{3}$ 时，杆切割磁感线的有效长度等于圆环半径，杆产生的感应电动势为 $E=Bav$，B错误；当 $\theta=0$ 时，回路的总电阻 $R_{1}=(2a+\pi a)R_{0}$，杆受的安培力 $F_{1}=BI_{1}l=B\cdot\frac{2Bav}{R_{1}}\cdot2a=\frac{4B^{2}av}{(\pi+2)R_{0}}$，C错误；当 $\theta=\frac{\pi}{3}$ 时，回路的总电阻 $R_{2}=(a+\frac{5}{3}\pi a)R_{0}$，杆受的安培力 $F_{2}=BI_{2}l'=B\cdot\frac{Bav}{R_{2}}\cdot a=\frac{3B^{2}av}{(5\pi+3)R_{0}}$，D正确。
}

\item
%% number: 21
%% paperName: 2012年四川理综第 21 题 6 分
%% typeId: 2
%% type: 多选题
%% chapter: 功和能
%% point: 动能定理
%% method: 无
%% score: 6
%% degree: 450
%% duplicateId: 0
%% body:
如图所示，劲度数为$k$的轻弹簧的一端固定在墙上，另一端与置于水平面上质量为$m$的物体接触（未连接），弹簧水平且无形变。用水平力$F$缓慢推动物体，在弹性限度内弹簧长度被压缩了$x_{0}$，此时物体静止。撤去$F$后，物体开始向左运动，运动的最大距离为$4x_{0}$。物体与水平面间的动摩擦因数为$\mu$，重力加速度为$g$。则\xzanswer{BD}
\onepicture{figs/21.png}
\fourchoices[answer=BD]
{撤去$F$后，物体先做匀加速运动，再做匀减速运动}
{撤去$F$后，物体刚运动时的加速度大小为$\frac{kx_{0}}{m}-\mu g$}
{物体做匀减速运动的时间为$2\sqrt{\frac{x_{0}}{\mu g}}$}
{物体开始向左运动到速度最大的过程中克服摩擦力做的功为$\mu mg(x_{0}-\frac{\mu mg}{k})$}
%% answer: BD
%% memo:
\memoanswer{
撤去F后的一段时间内，由于运动过程中弹力不断变化，物体先做变加速运动，后做变减速运动，再做匀减速运动，A错误；设撤去F后，物体刚运动时加速度为$a$，根据牛顿第二定律：$kx_0-\mu mg=ma$，解得$a=\frac{kx_0}{m}-\mu g$，B正确；物体做匀减速运动的位移为$3x_0$，由$3x_0=\frac{1}{2}\mu gt^2$，解得$t=\sqrt{\frac{6x_0}{\mu g}}$，C错误；当弹力与摩擦力大小相等时，速度最大，此时$kx_1=\mu mg$，该过程物体向左运动的位移为$x=x_0-x_1=x_0-\frac{\mu mg}{k}$，克服摩擦力做的功$W=\mu mg(x_0-\frac{\mu mg}{k})$，D正确。
}

\item
%% number: 22
%% paperName: 2012年四川理综第 22 题 6 分
%% typeId: 4
%% type: 实验题
%% chapter: 曲线运动
%% point: 研究平抛运动实验
%% method: 无
%% score: 6
%% degree: 650
%% duplicateId: 0
%% body:
某物理兴趣小组采用如图所示的装置深入研究平抛运动。质量分别为$m_{A}$和$m_{B}$的A、B小球处于同一高度，M为A球中心初始时在水平地面上的垂直投影。用小锤打击弹性金属片，使A球沿水平方向飞出，同时松开B球，B球自由下落。A球落到地面N点处，B球落到地面P点处。测得$m_{A}=0.04\Ukg$，$m_{B}=0.05\Ukg$，B球距地面的高度是$1.225\Um$，M、N点间的距离为$1.500\Um$，则B球落到P点的时间是\underline{\hspace{1.6cm}}$\Us$，A球落地时的动能是\underline{\hspace{1.6cm}}$\UJ$。（忽略空气阻力，$g$取$9.8\Umsq$）
\onepicture[align=right]{figs/22.png}
%% answer: 
\jdanswer{
0.5，0.66
}
%% memo:
\memoanswer{
B球做自由落体运动，有$h=\frac{1}{2}gt^{2}$，解得$t=0.5\Us$；A球做平抛运动，水平方向有$x=v_{0}t$，竖直方向有$h=\frac{1}{2}gt^{2}$，由动能定理，有$m_{A}gh=E_{k}-\frac{1}{2}m_{A}v_{0}^{2}$，解得$E_{k}=0.66\UJ$。
}

\item
%% number: 22
%% paperName: 2012年四川理综第 22 题 11 分
%% typeId: 4
%% type: 实验题
%% chapter: 电路
%% point: 伏安法测电阻
%% method: 无
%% score: 11
%% degree: 450
%% duplicateId: 0
%% body:
某学习小组的同学拟探究小灯泡L的伏安特性曲线，可供选用的器材如下：

小灯泡L，规格“$4.0\UV$，$0.7\UA$”；

电流表$A_{1}$，量程$3\UA$，内阻约为$0.1\UO$；

电流表$A_{2}$，量程$0.6\UA$，内阻$r_{2}=0.2\UO$；

电压表V，量程$3\UV$，内阻$r_{V}=0.9\UkO$；

标准电阻$R_{1}$，阻值$1\UO$；

标准电阻$R_{2}$，阻值$3\UkO$；

滑动变阻器$R$，阻值范围$0\sim10\UO$；

学生电源E，电动势$6\UV$，内阻不计；

开关S及导线若干。
\begin{enumerate}
	\item
	甲同学设计了如图1所示的电路来进行测量，当通过L的电流为$0.46\UA$时，电压表的示数如图2所示，此时L的电阻为\underline{\hspace{1.6cm}}$\UO$。
	\item
	乙同学又设计了如图3所示的电路来进行测量，电压表指针指在最大刻度时，加在L上的电压值是\underline{\hspace{1.6cm}}$\UV$。
	\item
	学习小组认为要想更准确地描绘出L完整的伏安特性曲线，需要重新设计电路。请你在乙同学的基础上利用所供器材，在图4所示的虚线框内补画出实验电路图，并在图上标明所选器材代号。
\end{enumerate}
\onepicture[h=2.6cm]{figs/22b.png}
%% answer: 
\jdanswer{
\begin{enumerate}
	\item
	$5\UO$
	\item
	$4\UV$
	\item
	见答图（电压表串联$R_{2}$改装，电流表$A_{2}$并联$R_{1}$改装后外接）
\end{enumerate}
}
%% memo:
\memoanswer{
①由图2可读出电压表的示数为$2.30\UV$，所以此时L的电阻为$2.30\UV/0.46\UA=5\UO$；

②电压表指针指在最大刻度时，电压表示数为$3.00\UV$，加在L上的电压值为$U_{L}=\frac{U_{V}}{r_{V}}(r_{V}+R_{2})=4\UV$；

③由于电压表量程太小，所以需要改装，将它与$R_{2}$串联即成为一个量程为$4.0\UV$的新的电压表；电流表$A_{1}$量程太大，不可用，可以将$A_{2}$改装：将它并联一个小电阻$R_{1}$则成为一个量程为$0.72\UA$的新的电流表了。由于灯泡的电阻远小于电压表内阻而与电流表内阻相差不多，属于小电阻，故用电流表外接法比较合适。答案如图1较合适，若用图2则不太合适。
\onepicture[h=2.8cm, align=right]{figs/22_an1.png}
\onepicture[h=2.8cm, align=right]{figs/22_an2.png}
}

\item
%% number: 23
%% paperName: 2012年四川理综第 23 题 16 分
%% typeId: 6
%% type: 计算题
%% chapter: 电路
%% point: 电功电功率
%% method: 无
%% score: 16
%% degree: 650
%% duplicateId: 0
%% body:
四川省“十二五”水利发展规划指出，若按现有供水能力测算，我省供水缺口极大，蓄引提水是目前解决供水问题的重要手段之一。某地要把河水抽高$20\Um$，进入蓄水池，用一台电动机通过传动效率为80\%的皮带，带动效率为60\%的离心水泵工作。工作电压为$380\UV$，此时输入电动机的电功率为$19\UkW$，电动机的内阻为$0.4\UO$。已知水的密度为$1\times10^{3}\Ukgmc$，重力加速度取$10\Umsq$。求：
\begin{enumerate}
	\item
	电动机内阻消耗的热功率；
	\item
	将蓄水池蓄入$864\Umc$的水需要的时间（不计进、出水口的水流速度）。
\end{enumerate}
%% answer: 
\jdanswer{
\begin{enumerate}
	\item
	$1\times10^{3}\UW$
	\item
	$2\times10^{4}\Us$
\end{enumerate}
}
%% memo:
\memoanswer{
(1)设电动机的电功率为$P$，则$P=UI$。设电动机内阻$r$上消耗的热功率为$P_{r}$，则$P_{r}=I^{2}r$，代入数据解得$P_{r}=1\times10^{3}\UW$。

(2)设蓄水总质量为$M$，所用抽水时间为$t$。已知抽水高度为$h$，容积为$V$，水的密度为$\rho$，则$M=\rho V$；设质量为$M$的河水增加的重力势能为$\Delta E_{p}$，则$\Delta E_{p}=Mgh$；设电动机的输出功率为$P_{0}$，则$P_{0}=P-P_{r}$；根据能量守恒定律得$P_{0}t\times60\%\times80\%=\Delta E_{p}$，代入数据解得$t=2\times10^{4}\Us$。
}

\item
%% number: 24
%% paperName: 2012年四川理综第 24 题 19 分
%% typeId: 6
%% type: 计算题
%% chapter: 电场
%% point: 带电粒子的直线运动
%% method: 无
%% score: 19
%% degree: 450
%% duplicateId: 0
%% body:
如图所示，ABCD为固定在竖直平面内的轨道，AB段光滑水平，BC段为光滑圆弧，对应的圆心角$\theta=37^\circ$，半径$r=2.5\Um$，CD段平直倾斜且粗糙，各段轨道均平滑连接，倾斜轨道所在区域有场强大小为$E=2\times10^{5}\UNC$、方向垂直于斜轨向下的匀强电场。质量$m=5\times10^{-2}\Ukg$、电荷量$q=+1\times10^{-6}\UC$的小物体（视为质点）被弹簧枪发射后，沿水平轨道向左滑行，在C点以速度$v_{0}=3\Ums$冲上斜轨。以小物体通过C点时为计时起点，$0.1\Us$以后，场强大小不变，方向反向。已知斜轨与小物体间的动摩擦因数$\mu=0.25$。设小物体的电荷量保持不变，取$g=10\Umsq$，$\sin37^\circ=0.6$，$\cos37^\circ=0.8$。
\begin{enumerate}
	\item
	求弹簧枪对小物体所做的功；
	\item
	在斜轨上小物体能到达的最高点为P，求CP的长度。
\end{enumerate}
\onepicture[align=right]{figs/24.png}
%% answer: 
\jdanswer{
\begin{enumerate}
	\item
	$0.475\UJ$
	\item
	$0.57\Um$
\end{enumerate}
}
%% memo:
\memoanswer{
(1)设弹簧枪对小物体做功为$W_{f}$，由动能定理得$W_{f}-mgr(1-\cos\theta)=\frac{1}{2}mv_{0}^{2}$，代入数据得$W_{f}=0.475\UJ$。

(2)取沿平直斜轨向上为正方向。设小物体通过C点进入电场后的加速度为$a_{1}$，由牛顿第二定律得$-mg\sin\theta-\mu(mg\cos\theta+qE)=ma_{1}$；小物体向上做匀减速运动，经$t_{1}=0.1\Us$后，速度达到$v_{1}$，有$v_{1}=v_{0}+a_{1}t_{1}$，联立以上方程可知$v_{1}=2.1\Ums$，设运动的位移为$s_{1}$，有$s_{1}=v_{0}t_{1}+\frac{1}{2}a_{1}t_{1}^{2}$；电场力反向后，设小物体的加速度为$a_{2}$，由牛顿第二定律得$-mg\sin\theta-\mu(mg\cos\theta-qE)=ma_{2}$；设小物体以此加速度运动到速度为0，运动的时间为$t_{2}$，位移为$s_{2}$，有$0=v_{1}+a_{2}t_{2}$，$s_{2}=v_{1}t_{2}+\frac{1}{2}a_{2}t_{2}^{2}$；设CP的长度为$s$，有$s=s_{1}+s_{2}$，联立相关方程，代入数据解得$s=0.57\Um$。
}

\item
%% number: 25
%% paperName: 2012年四川理综第 25 题 20 分
%% typeId: 6
%% type: 计算题
%% chapter: 磁场
%% point: 带电粒子在磁场中的运动
%% method: 无
%% score: 20
%% degree: 250
%% duplicateId: 0
%% body:
如图所示，水平虚线X下方区域分布着方向水平、垂直纸面向里、磁感应强度为$B$的匀强磁场，整个空间存在匀强电场（图中未画出）。质量为$m$、电荷量为$+q$的小球P静止于虚线X上方A点，在某一瞬间受到方向竖直向下、大小为$I$的冲量作用而做匀速直线运动。在A点右下方的磁场中有定点O，长为$l$的绝缘轻绳一端固定于O点，另一端连接不带电的、质量同为$m$的小球Q，自然下垂。保持轻绳伸直，向右拉起Q，直到绳与竖直方向有一小于$5^\circ$的夹角，在P开始运动的同时自由释放Q，Q到达O点正下方W点时的速率为$v_{0}$。P、Q两小球在W点发生正碰，碰后电场、磁场消失，两小球黏在一起运动。P、Q两小球均视为质点，P小球的电荷量保持不变，绳不可伸长，不计空气阻力，重力加速度为$g$。
\begin{enumerate}
	\item
	求匀强电场场强$E$的大小和P进入磁场时的速率$v$；
	\item
	若绳能承受的最大拉力为$F$，要使绳不断，$F$至少为多大？
	\item
	求A点距虚线X的距离$s$。
\end{enumerate}
\onepicture[align=right]{figs/25.png}
%% answer: 
\jdanswer{
\begin{enumerate}
	\item
	$E=\frac{mg}{q}$，$v=\frac{I}{m}$
	\item
	$\frac{(I+mv_0)^{2}}{2ml}+2mg$
	\item
	$(n+\frac{1}{4})\frac{2\pi I}{m}\sqrt{\frac{l}{g}}-\frac{\pi I}{2Bq}$（$n$为大于$\frac{m}{4Bq}\sqrt{\frac{g}{l}}-\frac{1}{4}$的整数）；或$(n+\frac{3}{4})\frac{2\pi I}{m}\sqrt{\frac{l}{g}}-\frac{\pi I}{2Bq}$（$n$为大于$\frac{m}{4Bq}\sqrt{\frac{g}{l}}-\frac{3}{4}$的整数）
\end{enumerate}
}
%% memo:
\memoanswer{
(1)设小球P所受电场力为$F_{1}$，则$F_{1}=qE$；在整个空间重力和电场力平衡，有$F_{1}=mg$，联立相关方程得$E=\frac{mg}{q}$。设小球P受到冲量后获得速度为$v$，由动量定理得$I=mv$，故$v=\frac{I}{m}$。

(2)设P、Q同向相碰后在W点的最大速度为$v_{m}$，由动量守恒定律得$mv+mv_{0}=(m+m)v_{m}$；此刻轻绳的张力也为最大，由牛顿运动定律得$F-(m+m)g=\frac{m+m}{l}v_{m}^{2}$，联立相关方程，得$F=\frac{(I+mv_{0})^{2}}{2ml}+2mg$。

(3)设P在X上方做匀速直线运动的时间为$t_{P1}$，则$t_{P1}=\frac{s}{v}$；设P在X下方做匀速圆周运动的时间为$t_{P2}$，则$t_{P2}=\frac{\pi m}{2Bq}$。设小球Q从开始运动到与P球反向相碰的运动时间为$t_{Q}$，由单摆周期性，有$t_{Q}=(n+\frac{1}{4})2\pi\sqrt{\frac{l}{g}}$，由题意，有$t_{Q}=t_{P1}+t_{P2}$，联立相关方程，得$s=(n+\frac{1}{4})\frac{2\pi I}{m}\sqrt{\frac{l}{g}}-\frac{\pi I}{2Bq}$，[$n$为大于$\frac{m}{4Bq}\sqrt{\frac{g}{l}}-\frac{1}{4}$的整数]。若小球Q从开始运动到与P球同向相碰，由单摆周期性，有$t_{Q}=(n+\frac{3}{4})2\pi\sqrt{\frac{l}{g}}$，同理可得$s=(n+\frac{3}{4})\frac{2\pi I}{m}\sqrt{\frac{l}{g}}-\frac{\pi I}{2Bq}$，[$n$为大于$\frac{m}{4Bq}\sqrt{\frac{g}{l}}-\frac{3}{4}$的整数]。
}

\end{enumerate}

\end{document}
